\documentclass[letterpaper, openright]{report}
\usepackage{titlesec}
\usepackage[english]{babel}
\usepackage[utf8x]{inputenc}
% package for including graphics with figure-environment
\usepackage{graphicx}
\usepackage{hyperref}
\usepackage{amsmath}
\usepackage{csquotes}% Recommended
% colors for hyperlinks
% colored borders (false) colored text (true)
\hypersetup{colorlinks=true,citecolor=black,filecolor=black,linkcolor=black,urlcolor=black}
% acronyms
\usepackage{acronym} 
%for inserting pdf files
\usepackage[final]{pdfpages}
%bibliography
%\usepackage{cite}
\usepackage{natbib}
%\usepackage{biblatex}
\usepackage[nottoc]{tocbibind}%needed for displaying bib and in table of contents
%Wordcount
\usepackage{bashful}
\pagenumbering{roman}

% package for header
\usepackage[automark]{scrpage2}
\pagestyle{scrheadings}
\ihead[]{Guillaume Auvray}
\ohead[]{\today}
\cfoot[]{\pagemark} 
\setheadsepline[122mm]{0.3mm}
\begin{document}
	\title{
	\begin{figure}[!ht]
		\flushleft
        	\centering
			\includegraphics[width=1.0 \textwidth]{IMG/LogoPulseW.png}
	\end{figure}
    \vspace{1cm}
    \begingroup
    \leftskip=0cm plus 0.5fil \rightskip=0cm plus -0.5fil \parfillskip=0cm plus 1fil
	\Large \textbf{Binaural Spatialised Music in \ac{CGI} and non-\ac{CGI} Media : \newline  \textit{Spatial Electroacoustic composition as part of fully immersive and interactive environments as \ac{VR} and other 360° Media.}} \\ 
    \endgroup
    \vspace{1cm}
    \Large \textit{Master Thesis}
	}
	
	
	% if you are the only author, you might use the following
	% \author{Name of student}
	
	% Insert here your name and correct mail address
	\author{\Large \href{mailto:guillaumeauvray@hotmail.com}{Guillaume Auvray}
	\vspace{0.1cm}}
            \newcommand{\detailtexcount}{%
  \immediate\write18{texcount -merge -sum -incbib -dir \jobname.tex > \jobname.wcdetail }%
  \verbatiminput{\jobname.wcdetail}%
}

\newcommand{\quickwordcount}{%
  \immediate\write18{texcount -1 -sum -merge \jobname.tex > \jobname-words.sum }%
  \input{\jobname-words.sum}%
}

\newcommand{\quickcharcount}{%
  \immediate\write18{texcount -1 -sum -merge -char \jobname.tex > \jobname-chars.sum }%
  \input{\jobname-chars.sum} characters (not including spaces)%
}
	
	% name of the course and module
	\date{
	\large Course: MA Scoring for Visual Media and Film \\ 
	\vspace{0.1cm}
	\large Supervisor: Dr. Natasa Paterson-Paulberg \\
    \large Advisor: Greg Clooney \\
    \vspace{0.3cm}
    \large Word count : \textbf{\quickwordcount} \\
    \vspace{0.1cm}
    %\today
    \begin{figure}[h!]
    	\centering
    	\includegraphics[width=0.45\textwidth]{IMG/BannerDIT.jpg}
	\end{figure}
	}
	\maketitle 
	\newpage

    \begin{abstract}
\textit{Abstract to be written...}
	\end{abstract}

\addcontentsline{toc}{part}{Signed declaration}
\begin{center}

  \huge \textbf{Signed Declaration}
  
\end{center} 
\vspace{0.3cm}
I declare that this thesis has been composed solely by myself and that it has not been submitted, in whole or in part, in any previous application for a degree. Except where states otherwise by reference or acknowledgment, the work presented is entirely my own.

\vspace{12cm}

\begin{flushright}

  \large Date and Signature
  
\end{flushright}

\newpage

\addcontentsline{toc}{part}{Dedication}
\begin{center}

  \huge \textbf{Dedication}
  
\end{center} 
\vspace{0.3cm}
This thesis is dedicated to... I dint know yet :D, Well, actually I do, but I'm too lazy to write it out now. I will do it at the end.
 

	\newpage   
	\thispagestyle{empty}
	
   	\listoffigures
	\listoftables
    \newpage
	
    \title{\textbf{List of Acronym}}   
    \input{acro_list}
 	\newpage   
    \newpage
	\tableofcontents
	\newpage



\chapter{Introduction} % (fold)
\label{sec:introduction}
\pagenumbering{arabic}

	\section{Introduction} % (fold)
	\label{sub:intro-intro}
\textit{intro-Intro has to be written...}	
% subsection intro-intro (end)
    
    \section{Project Description} % (fold)
	\label{sub:instro-projdesc}
\textit{intro-projdesc has to be written...}	
% subsection intro-projdesc (end)

			\subsection{Technology} % (fold)
			\label{subsub:intro-projdesc-tech}       
\textit{intro-projdesc-tech has to be written...}	
% subsubsection intro-projdesc-tech (end)

			\subsection{Composition} % (fold)
			\label{subsub:intro-projdesc-comp}
\textit{intro-projdesc-comp has to be written...}
% subsubsection intro-projdesc-comp (end)

			\subsection{\ac{VR} Experience} % (fold)
			\label{subsub:intro-projdesc-VR_Exp}
\textit{intro-projdesc-VRExp has to be written...}	
			% sub-subsection intro-projdesc-VR_Exp (end)
            
	% subsection introduction (end)

    \section{Genre} % (fold)
	\label{sub:genre}
\textit{genre has to be written...}
			% sub-subsection intro-genre (end)

			\subsection{Interactive Media} % (fold)
			\label{subsub:intro-genre-iMedia}
\textit{intro-genre-iMedia has to be written...}
			% sub-subsection intro-genre-iMedia (end)

			\subsection{\ac{XR}} % (fold)
			\label{subsub:intro-genre-XR}
\textit{intro-genre-XR has to be written...}
			% sub-subsection intro-genre-XR (end)
            
            
	% subsection genre (end)
    
    \section{Goals and Motivation} % (fold)
	\label{sub:intro-motivation}
\textit{intro-motivation has to be written...}
	% subsection intro-motivation (end)
    
    \section{Key Contributions} % (fold)
	\label{sub:intro-contribution}
\textit{intro-contribution has to be written...}
	% subsection intro-contribution (end)

    \section{Thesis Outline} % (fold)
	\label{sub:intro-outline}
\textit{intro-outline has to be written...}
	% subsection intro-outline (end)

    \section{Conclusion} % (fold)
	\label{sub:intro-conclusion}
This chapter covered a few of the numerous preambles fueling the concerns we have to bear in mind in order to understand the areas we will have to cover. We took an overall look at the technological aspects which are involved in the work, the compositional background and the interactive field we will have to deal with. \newline
The next chapters will cover the Background of the research, followed by the Review and Methodology. Thereafter, we will explain and have an in-depth look at the compositional work related to the research.
	% subsection intro-conclusion (end)
    
% section introduction (end)

	\vspace{0.5cm}


\newpage
	\newpage
\chapter{Background} % (fold)
\label{sec:background}

	\section{Introduction} % (fold)
	\label{sub:back-intro}
In this chapter, we will discuss several points which are directly related to the perception of the music and how we can maintain the attention of the listener at each place of the orchestra. These elements will form the biggest part of the background we will build our research on.
	% subsection back-intro (end)
    
    \section{Spatial Music} % (fold)
	\label{sub:back-spatialM}

			\subsection{Enda Bates} % (fold)
			\label{subsub:back-spatialM-eBates}
\textit{back-spatialM-eBates has to be written...}
			% sub-subsection back-spatialM-eBates (end)

			\subsection{Varese} % (fold)
			\label{subsub:back-spatialM-Varese}
\textit{back-spatialM-Varese has to be written...}
			% sub-subsection back-spatialM-Varese (end)    
            
	% subsection back-spatialM (end)

    \section{Musical Style} % (fold)
	\label{sub:back-style}
    
			\subsection{Spatial Music} % (fold)
			\label{subsub:back-style-spatialM}
\textit{back-style-spatial has to be written...}
			% sub-subsection back-style-spatialM (end)
            
            \subsection{Music storytelling: Opera} % (fold)
			\label{subsub:back-style-opera}
\textit{back-style-spatial has to be written...}
			% sub-subsection back-style-opera (end)
            
			\subsection{Ligeti} % (fold)
			\label{subsub:back-style-ligeti}
\textit{back-style-ligeti has to be written...}
			% sub-subsection back-style-ligeti (end)   
 
			\subsection{Bjork} % (fold)
			\label{subsub:back-style-Bjork}
\textit{back-style-Bjork has to be written...}
			% sub-subsection back-style-Bjork (end)  
 
			\subsection{Electro-acoustic or Tape music} % (fold)
			\label{subsub:back-style-electro}
\textit{back-style-electro has to be written...}
			% sub-subsection back-style-electro (end)                        
	% subsection back-style (end)
    
    \section{Immersive Audio Technologies} % (fold)
	\label{sub:back-IAT}
\textit{back-IAT has to be written...}

			\subsection{Immersive Media} % (fold)
			\label{subsub:back-IAT-iMedia}
\textit{back-IAT-iMedia has to be written...}
 			% sub-subsection back-IAT-iMedia (end)

			\subsection{Higher Order Ambisonics} % (fold)
			\label{subsub:back-IAT-HOA}
\textit{back-IAT-HOA has to be written...}
 			% sub-subsection back-IAT-iMedia (end)    

	% subsection back-IAT (end)
    
    \section{Conclusion} % (fold)
	\label{sub:back-conclusion}
This chapter covered the background needed in order to achieve properly the research that has to be made.  We discussed several points which are directly related to the perception of the music and how we can maintain the attention of the listener at each place of the orchestra.\newline
The next chapters will cover the Review of works and literature already existing in the domain which has to be studied.
	% subsection back-conclusion (end)
% section background (end)

	\vspace{0.5cm}
   

\newpage
	\newpage
\chapter{Review} % (fold)
\label{sec:review}

	\section{Introduction} % (fold)
	\label{sub:rev-intro}
In this chapter we will be reviewing some important works or literature written in relation to the background knowledge we have to build in order to achieve a proper understanding of the research.
	% subsection rev-intro (end)
    
    \section{Spatial Music} % (fold)
	\label{sub:rev-spatialM}

			\subsection{Enda Bates} % (fold)
			\label{subsub:rev-spatialM-eBates}
\textit{rev-spatialM-eBates has to be written...}
			% sub-subsection rev-spatialM-eBates (end)

			\subsection{Varese} % (fold)
			\label{subsub:rev-spatialM-Varese}
\textit{rev-spatialM-Varese has to be written...}
			% sub-subsection rev-spatialM-Varese (end)    
            
	% subsection rev-spatialM (end)

    \section{Musical Style} % (fold)
	\label{sub:rev-style}
    
			\subsection{Spatial Music} % (fold)
			\label{subsub:rev-style-spatialM}
\textit{rev-style-spatialM has to be written...}
			% sub-subsection rev-style-spatialM (end)
            
            \subsection{Music storytelling: Opera} % (fold)
			\label{subsub:rev-style-opera}
\textit{rev-style-spatial has to be written...}
			% sub-subsection rev-style-opera (end)
            
			\subsection{Ligeti} % (fold)
			\label{subsub:rev-style-ligeti}
\textit{rev-style-ligeti has to be written...}
			% sub-subsection rev-style-ligeti (end)   
 
			\subsection{Bjork} % (fold)
			\label{subsub:rev-style-Bjork}
\textit{rev-style-Bjork has to be written...}
			% sub-subsection rev-style-Bjork (end)  
 
			\subsection{Electro-acoustic or Tape music} % (fold)
			\label{subsub:rev-style-electro}
\textit{rev-style-electro has to be written...}
			% sub-subsection rev-style-electro (end)                        
	% subsection rev-style (end)
    
    \section{Immersive Audio Technologies} % (fold)
	\label{sub:rev-IAT}
\textit{rev-IAT has to be written...}

			\subsection{Immersive Media} % (fold)
			\label{subsub:rev-IAT-iMedia}
\textit{rev-IAT-iMedia has to be written...}
 			% sub-subsection rev-IAT-iMedia (end)

			\subsection{Higher Order Ambisonics} % (fold)
			\label{subsub:rev-IAT-HOA}
\textit{rev-IAT-HOA has to be written...}
 			% sub-subsection rev-IAT-iMedia (end)    

	% subsection rev-IAT (end)   
    
    \section{Conclusion} % (fold)
	\label{sub:rev-conclusion}
In this chapter we will review all the background knowledge needed in order to build the understanding the research implies.  We discussed several points which are directly related to the perception of the music and how we can maintain the attention of the listener at each place of the orchestra.\newline
The next chapters will cover the Methodology that has to be followed in order to achieve the final result of the work.
	% subsection rev-conclusion (end)    
% section review (end)

	\vspace{0.5cm}

	\newpage

\chapter{Methodology} % (fold)
\label{sec:methodo}

	\section{Introduction} % (fold)
	\label{sub:methodo-intro}
In this chapter we will explain how the work was done. All the thoughts and decisions undertaken during the elaboration will be listed in three major sections: 
\begin{enumerate}
\item Implementation
\item Composition
\item Technical
\end{enumerate}
	% subsection intro (end)
    
    \section{Implementation} % (fold)
	\label{sub:methodo-implement}

			\subsection{Musicians} % (fold)
			\label{subsub:methodo-implement-music}
\textit{methodo-implement-music has to be written...}
			% sub-subsection methodo-implement-music (end) 
            
            \subsection{Location} % (fold)
			\label{subsub:methodo-implement-location}
\textit{methodo-implement-location has to be written...}
			% sub-subsection methodo-implement-location (end) 
            
            \subsection{Recording Rigs} % (fold)
			\label{subsub:methodo-implement-rigs}
\textit{methodo-implement-rigs has to be written...}
			% sub-subsection methodo-implement-rigs (end) 
      % subsection methodo-implement (end)

    \section{Post-Production} % (fold)
	\label{sub:methodo-postprod} 

            \subsection{Editing} % (fold)
			\label{subsub:methodo-postprod-edit}
\textit{methodo-postprod-edit has to be written...}
			% sub-subsection methodo-postprod-edit (end) 

				\subsubsection{Audio} % (fold)
				\label{subsub:methodo-postprod-edit-audio}
\textit{methodo-postprod-edit-audio has to be written...}
				% sub-subsection methodo-postprod-audio (end)
            
            	\subsubsection{Video} % (fold)
				\label{subsub:methodo-postprod-video}
\textit{methodo-postprod-edit-video has to be written...}
				% sub-subsection methodo-postprod-video (end) 
            
            \subsection{\ac{VFX}} % (fold)
			\label{subsub:methodo-postprod-vfx}
\textit{methodo-postprod-vfx has to be written...}
			% sub-subsection methodo-postprod-vfx (end)

		% subsection methodo-postprod (end)
    
        \section{Rendering} % (fold)
		\label{subsub:methodo-render}
\textit{methodo-render has to be written...}
		% sub-subsection methodo-render (end)

    
    \section{Conclusion} % (fold)
	\label{sub:conclusion}
We will face a couple of important choices during the recordings and the implementation. Several versions will be made and placed on private 360 Youtube videos in order to discuss the final results with musicians, non-musicians, and technicians before rendering the final application or video.
\newline
In the beginning we consider working with Google Resonance. Google Resonance will be used in a second stage. Google Resonance will be powerful and very effective when we will be in a full \ac{CG} environment. Each musician will be 3D scanned and placed in the \ac{VR} environment. Each 3D musician will be live animated with Motion-Data sensors clipped on their desks and linked to the skeleton of the 3D mesh inside the Game-Engine. Google Resonance will be fed with all the microphone tracks and will render a binaural track in function of where the viewer will situate himself in space. At that level we will not have to struggle with that amount of cameras anymore. Because the environment will be entirely computer generated, the viewer will have entire freedom to move virtually on and around the stage and perceive the sound as if he would be there.
\newline
This technology is growing very fast, but as soon as we can get it on an orchestra or a television news presenter, it will be soon available for sport and other entertainment companies. Their spectators will wear their headset and they will be able to be transported to the stadium and evolve on the court as if they would walk around with the players. As we have seen in \citep{Alvarez2018}, the digital human "Siren" at the GDC2018 the collaboration between Epic Games, Vicon, Cubic Motion, 3Lateral and Tencent. This amazing interaction between motion data, CGI and a game engine will bring us towards the future of broadcasting entertainment. We don't have to forget the power of sound and music, which makes 70\% of the immersive reality. This is where we as technologists, composers and musicians have to jump into.
	% subsection conclusion (end)  
% section methodo (end)

	\vspace{0.5cm}

	\newpage
 
\chapter{The work} % (fold)
\label{sec:work}

	\section{Introduction} % (fold)
	\label{sub:work-intro}
Initially the work was intended to be made on a big ensemble, like a chambre orchestra. But we decided to think smaller in order to go in-depth with a string quartet. The understanding we will build with this small instrumentation will be helpful to develop it later for orchestral ensembles.
	% subsection work-intro (end)

    \section{Musical Analysis} % (fold)
	\label{sub:work-analyse}
As soon as the score will be written, an in-depth analysis of the choice made and the musical gesture will be done and divided in these three categories : 
\begin{enumerate}
\item Structure and form
\item Counterpoint and Melodic Material
\item Harmonic material
\end{enumerate}

            \subsection{Structure and Form} % (fold)
			\label{subsub:work-analyse-struct}
\textit{work-analyse-struct has to be written...}
			% sub-subsection work-analyse-struct (end)

            \subsection{Rhytmic Material} % (fold)
			\label{subsub:work-analyse-rythm}
\textit{work-analyse-rythm has to be written...}
			% sub-subsection work-analyse-rythm (end)

            \subsection{Melodic Material} % (fold)
			\label{subsub:work-analyse-melo}
\textit{work-analyse-melo has to be written...}
			% sub-subsection work-analyse-melo (end)

	% subsection work-analyse (end)
    
    \section{Immersive Behaviours} % (fold)
	\label{sub:work-immersive}
In this section we will analyse the relation between the visual and acoustic aspect of the work. How does the viewer perceive the music? At which point and how would we as composers, be able to enhance a reaction by the viewer? Which elements are maintaining the interest of the  viewer in en \ac{VR} environment? All these questions will have to be answered and analysed.
	% subsection ch3 (end)
    
            \subsection{Psychological aspects} % (fold)
			\label{subsub:work-immersive-psycho}
\textit{work-immersive-psycho has to be written...}
			% sub-subsection work-immersive-psycho (end)
    
    \section{Conclusion} % (fold)
	\label{sub:work-conclusion}
\textit{work-conclusion has to be written...}
	% subsection conclusion (end)
% section work (end)

	\vspace{0.5cm}    

	\newpage
    
\chapter{Discussion} % (fold)
\label{sec:discussion}

	\section{Introduction} % (fold)
	\label{sub:disc-intro}
In this chapter we will discuss some elements which are essential in order to achieve the expected results. 
	% subsection disc-intro (end)
    
    \section{Critique, Evaluation and Preparation} % (fold)
	\label{sub:disc-crit}
\textit{disc-crit has to be written...}
	% subsection disc-crit (end)    
% section discussion (end)

	\newpage
    
\chapter{Conclusion} % (fold)
\label{sec:conclusion}

	\section{Introduction} % (fold)
	\label{sub:concl-intro}
\textit{concl-intro has to be written...}
	% subsection concl-intro (end)
    
	\section{Summary of Thesis} % (fold)
	\label{sub:concl-summary}
\textit{concl-summary has to be written...}
	% subsection concl-summary (end)
    
	\section{Major Findings and Answers} % (fold)
	\label{sub:concl-answers}
\textit{concl-answers has to be written...}
	% subsection concl-answers (end)
    
	\section{Future Works} % (fold)
	\label{sub:concl-future}
\textit{concl-future has to be written...}
	% subsection concl-future (end)    

	\section{Conclusion to the Thesis} % (fold)
	\label{sub:concl-finalconcl}
\textit{concl-finalconcl has to be written...}
	% subsection concl-finalconcl (end)   

	\vspace{0.5cm}       

	\newpage

% The bibliography
\bibliographystyle{agsm}
	\bibliography{biblio} % expects file "biblio.bib"
% section bib (end)

\vspace{0.5cm}  

\addcontentsline{toc}{part}{Appendix}



	\newpage
    
\end{document}